\documentclass[12pt, a4paper]{article}

% ── Packages ──────────────────────────────────────────────────────────────────
\usepackage[margin=2.5cm]{geometry}
\usepackage{booktabs}           % pretty tables
\usepackage{tabularx}           % flexible table widths
\usepackage{caption}            % caption formatting
\usepackage{graphicx}           % figures
\usepackage{amsmath}            % math equations
\usepackage{amssymb}
\usepackage{float}              % figure/table placement [H]
\usepackage{hyperref}           % clickable links
\usepackage{natbib}             % citations
\usepackage{setspace}           % line spacing
\usepackage{parskip}            % space between paragraphs, no indent
\usepackage{microtype}          % better typography
\usepackage{xcolor}
\usepackage{titlesec}           % section formatting
\usepackage{lmodern}            % better font rendering
\usepackage{threeparttable}     % table notes
\usepackage{siunitx}            % number formatting in tables

% ── Hyperlink style ───────────────────────────────────────────────────────────
\hypersetup{
    colorlinks=true,
    linkcolor=black,
    urlcolor=blue,
    citecolor=black
}

% ── Section formatting ────────────────────────────────────────────────────────
\titleformat{\section}{\large\bfseries}{\thesection.}{0.5em}{}
\titleformat{\subsection}{\normalsize\bfseries}{\thesubsection.}{0.5em}{}

% ── Line spacing ──────────────────────────────────────────────────────────────
\onehalfspacing

% ── Title info ────────────────────────────────────────────────────────────────
\title{
    \vspace{-1cm}
    {\large DA4 Assignment 2} \\[0.4em]
    {\LARGE \textbf{CO2 Emissions and Economic Activity: \\A Panel Data Analysis}}
}
\author{Mai Ngoc Nong}
\date{\today}

% ══════════════════════════════════════════════════════════════════════════════
\begin{document}
% ══════════════════════════════════════════════════════════════════════════════

\maketitle
\thispagestyle{empty}

\vspace{1em}
\noindent\textbf{Data and code:} \href{https://github.com/MiaHenstridge/panel-regression-co2-emission-gdp-causal-relationship.git}{github.com/MiaHenstridge/panel-regression-co2-emission-gdp-causal-relationship.git}


% ══════════════════════════════════════════════════════════════════════════════
\section{Introduction and Data Description}
% ══════════════════════════════════════════════════════════════════════════════

This analysis uses two series from the World Bank World Development Indicators: CO2 emissions per capita (metric tons, indicator \texttt{EN.GHG.CO2.PC.CE.AR5}) and GDP per capita in PPP current international dollars (\texttt{NY.GDP.PCAP.PP.CD}), covering 1992 to 2024, to explore the effect of economic activities on CO2 emissions. After dropping countries with low coverage (below 50\% non-missing observations), the dataset contains 192 units. Additionally, the analysis also uses Urbanization population as percentage of total population (\texttt{SP.URB.TOTL.IN.ZS}) as a confounder variable. CO2 emissions per capita and GDP per capita variables are log-transformed in all regressions, so their coefficients are interpretable as elasticities. Since urban population is already measured in percentage point, no further transformation is necessary.

\subsection{Coverage and Sample Selection}

% ── Table 1: Coverage summary ─────────────────────────────────────────────────
\begin{table}[H]
\centering
\begin{threeparttable}
\caption{Sample Coverage}
\label{tab:coverage}
\begin{tabularx}{0.6\textwidth}{Xr}
\toprule
Initial countries   & 217 \\
Dropped ($< 50\%$ coverage) & 25 \\
Final sample        & 192 \\
Years               & 1992--2024 \\
Total observations  & 6336 \\
\bottomrule
\end{tabularx}
\begin{tablenotes}
\small
\item Countries with fewer than 50\% of years non-missing for both CO2 and GDP are dropped.
\end{tablenotes}
\end{threeparttable}
\end{table}

\subsection{Descriptive Statistics}

% ── Table 2: Summary stats ────────────────────────────────────────────────────
\begin{table}[H]
\centering
\begin{threeparttable}
\caption{Summary Statistics}
\label{tab:sumstats}
\begin{tabularx}{\textwidth}{Xrrrrr}
\toprule
\textbf{Variable} & \textbf{Mean} & \textbf{SD} & \textbf{Min} & \textbf{Max} & \textbf{Missing(\%)}\\
\midrule
CO2 per capita (metric tons)        & 4.92      & 9.59      & 0.00      & 202.87    & 0.00\\[0.3em]
GDP per capita (PPP, current USD)   & 17683.99  & 21614.46  & 254.39    & 180939.44 & 1.74\\[0.3em]
Urban population (\%)               & 56.73     & 23.77     & 5.66      & 100.00    & 0.00\\[0.3em]
\bottomrule
\end{tabularx}
\end{threeparttable}
\end{table}

% ══════════════════════════════════════════════════════════════════════════════
\section{Empirical Strategy}
% ══════════════════════════════════════════════════════════════════════════════

All models are estimated with the log-transformed outcome variable and causal variable, yielding elasticity interpretations. Let $y_{it}$ denote log CO2 per capita and $x_{it}$ log GDP per capita for country $i$ in year $t$.

\subsection{Cross-Sectional OLS}

\begin{equation}
y_{it} = \alpha + \beta \, x_{it} + \varepsilon_{it}
\label{eq:cross}
\end{equation}

Estimated separately for 2005 and the most recent available year. HC3 heteroskedasticity-robust standard errors are used throughout.

\subsection{First Difference Model with Time Trend and Lags}

\begin{equation}
\Delta y_{it} = \alpha + \sum_t \gamma_t \mathbf{1}[year = t] +  \beta_0 \, \Delta x_{it} + \sum_{k=1}^{K} \beta_k \, \Delta x_{i,t-k} + \varepsilon_{it}
\label{eq:fd}
\end{equation}

where $\Delta y_{it} = y_{it} - y_{i,t-1}$, $\Delta x_{it} = x_{it} - x_{i,t-1}$,  $K \in \{0, 2, 6\}$. First differencing removes any time-invariant country-level confounders. $\sum_t \gamma_t \mathbf{1}[year = t]$ denote year dummies absorbing global shocks common to all countries in a given year. Standard errors are clusted at the country level.

\subsection{Fixed Effects Model with Time and Country Fixed Effects}

\begin{equation}
y_{it} = \alpha_i + \sum_t \gamma_t \mathbf{1}[year = t] + \beta \, x_{it} + \varepsilon_{it}
\label{eq:fe}
\end{equation}

where $\alpha_i$ are country fixed effects.

% ══════════════════════════════════════════════════════════════════════════════
\section{Results}
% ══════════════════════════════════════════════════════════════════════════════

The results of the estimated OLS, First Difference, and Fixed Effects regression models are summarized in table \ref{tab:results}. The cross-sectional OLS estimates suggest a positive association between GDP per capita and CO2 emissions per capita. In 2005, a 1\% higher GDP per capita across countries is associated with a 1.22\% higher CO2 per capita, falling slightly to 1.06\% in the most recent year, suggesting a modest decoupling of emissions from income over time. However, these models' estimates are likely upward biased as they do not account for country-specific differences.

The first difference and fix effects models show a dampened estimates after absorbing time-invariant counfounders. The FD models which focuses on within-country year-on-year variation, yield substantially smaller elasticities: a 1\% increase in GDP per capita growth is associated with 0.27\% - 0.32\% increase in CO2 per capita growth in the same year across the lag specifications. The lagged GDP per capita growth terms are largely insignificant across $K=2$ and $K=6$, suggesting that the effect of GDP per capita growth on emissions materializes quickly rather than exhibiting over the long term.

Finally, the FE model which exploits the within-country variation over time while controlling for country-specific levels and global annual shocks, yields an elasticity of 0.44\%. This means a 1\% higher in GDP per capita growth compared to a country's average is associated with a 0.44\% increase in CO2 emission per capita growth compared to the same country's average.

The consistent decline in the elasticity observed from the cross-section OLS to FD and FE suggests that a substantial portion of the raw GDP-CO2 correlation is likely attributable to between-country differences rather than a causal within-country effect of economic growth on emissions.

% ── Table 3: Main results ─────────────────────────────────────────────────────
\begin{table}[H]
\centering
\begin{threeparttable}
\caption{Regression Results: CO2-GDP Elasticity}
\label{tab:results}
\begin{tabularx}{\textwidth}{Xcccccc}
\toprule
& \multicolumn{2}{c}{\textbf{Cross-section OLS}} & \multicolumn{3}{c}{\textbf{FD}} & \textbf{FE} \\
\cmidrule(lr){2-3} \cmidrule(lr){4-6} \cmidrule(lr){7-7}
& 2005 & Latest & $K=0$ & $K=2$ & $K=6$ & \\
\midrule
log GDP p.c.                    & 1.2180*** & 1.0599*** & - & - & - & 0.4414*** \\
                                & (0.0679)  & (0.0646)  & - & - & - & (0.0747) \\[0.3em]

$\Delta$ log GDP p.c.           & - & - & 0.3195*** & 0.2708*** & 0.2687*** & - \\
                                & - & - & (0.0481)  & (0.0482)  & (0.0520)  & - \\[0.3em]

$\Delta$ log GDP p.c. $_{t-1}$  & - & - & - & 0.0191 & 0.0381 & -\\
                                & - & - & - & (0.0472) & (0.0495) & -\\[0.3em]

$\Delta$ log GDP p.c. $_{t-2}$  & - & - & - & 0.0516** & 0.0202 & -\\
                                & - & - & - & (0.0214) & (0.0246) & -\\[0.3em]

$\Delta$ log GDP p.c. $_{t-3}$  & - & - & - & - & 0.0232 & -\\
                                & - & - & - & - & (0.0344) & -\\[0.3em]

$\Delta$ log GDP p.c. $_{t-4}$  & - & - & - & - & 0.0410 & -\\
                                & - & - & - & - & (0.0342) & -\\[0.3em]

$\Delta$ log GDP p.c. $_{t-5}$  & - & - & - & - & -0.0366 & -\\
                                & - & - & - & - & (0.0370) & -\\[0.3em] 

$\Delta$ log GDP p.c. $_{t-6}$  & - & - & - & - & 0.0457 & -\\
                                & - & - & - & - & (0.0322) & -\\[0.3em]                                                              
\midrule
Time FE                 & No & No & Yes & Yes & Yes & Yes \\
Country FE              & No & No & No  & No  & No  & Yes \\
Observations            & 187   & 179   & 5971  & 5591  & 4831  & 6161   \\
R$^2$                   & 0.640 & 0.634 & 0.023 & 0.019 & 0.019 & 0.069 \\
\bottomrule
\end{tabularx}
\begin{tablenotes}
\small
\item Standard errors in parentheses. HC3 robust SEs are reported for OLS and FD models; clustered SEs are reported for FD and FE models. R$^2$ within are reported for FD and FE models. Due to log transformation, models are estimated using only observations with both positive CO2 emission per capita and GDP per capita. $^{*}p<0.10$, $^{**}p<0.05$, $^{***}p<0.01$.
\end{tablenotes}
\end{threeparttable}
\end{table}

% ══════════════════════════════════════════════════════════════════════════════
\section{Extension of analysis: Urban Population as a Confounder}
% ══════════════════════════════════════════════════════════════════════════════

To examine whether the GDP-CO2 relationship is confounded by country-specific attributes, this analysis adds countries' urban population (\%) as a potential confounder to models 1, 4, and 6. Urbanization is expected to be positively correlated with CO2 emissions independently of GDP per capita, driven by increased energy demand for transport, residential heating and cooling, and industrial activities concentrated in cities. Thus, a country with a higher urban population is likely to have higher emission intensity and omitting this variable may have biased the GDP coefficient upward.

Table \ref{tab:results-confounder} reports the results after adding urban population as a confounder to models 1, 4, and 6. The GDP coefficients are dampened across the cross-section OLS and FE models compared to the initial estimates in Table \ref{tab:results}, falling from 1.22 to 1.18 in the cross-section OLS and from 0.44 to 0.41 in the FE model. This result is consistent with the hypothesis that part of the baseline GDP coefficient was inflated by omitted variable bias. The dampening is however modest, suggesting that economic activities retains most of its explanatory power independently of urbanization.

The urbanization coefficient is insignificant in the OLS model but positive and significant in the FD and FE models. In particular, the FE estimate implies that a one percentage point increase in a country's urban population compared to its average is associated with a 1.9\% increase compared to its average in CO2 emission per capita. The FD coefficient implies that a one percentage point increase in year-on-year growth of urban population is associated with a 1.1\% increase in year-on-year CO2 emission per capita.

% ── Table 4: Results with Confounder ─────────────────────────────────────────────────────
\begin{table}[H]
\centering
\begin{threeparttable}
\caption{Regression Results: CO2-GDP Elasticity with Industry Share as Confounder}
\label{tab:results-confounder}
\begin{tabularx}{\textwidth}{Xccc}
\toprule
& Cross-section OLS (2005) & FD ($K=2$) & FE \\
\midrule
log GDP p.c.            & 1.1822***  & - & 0.4136*** \\
                        & (0.1148)   & - & (0.0695) \\[0.3em]

$\Delta$ log GDP p.c.   & - & 0.2701*** & - \\
                        & - & (0.0481)  & - \\[0.3em]

$\Delta$ log GDP p.c. $_{t-1}$  & - & 0.0173   & -\\
                                & - & (0.0472) & -\\[0.3em]

$\Delta$ log GDP p.c. $_{t-2}$  & - & 0.0483** & -\\
                                & - & (0.0212) & -\\[0.3em] 
                                
urban population (\%)             & 0.0024 & - & 0.0187*** \\
                                & (0.0076) & - & (0.0048)  \\[0.3em]

$\Delta$ urban population (\%)    & - & 0.0113***   & - \\
                                & - & (0.0037) & - \\[0.3em]
\midrule
Time FE                 & No & Yes & Yes \\
Country FE              & No & No  & Yes \\
Observations            & 187   & 5591   & 6161 \\
R$^2$                   & 0.640 & 0.021 & 0.100 \\
\bottomrule
\end{tabularx}
\begin{tablenotes}
\small
\item Standard errors in parentheses. HC3 robust SEs are reported for OLS and FD models; clustered SEs are reported for FD and FE models. R$^2$ within are reported for FD and FE models. Due to log transformation, models are estimated using only observations with both positive CO2 emission per capita and GDP per capita. $^{*}p<0.10$, $^{**}p<0.05$, $^{***}p<0.01$.
\end{tablenotes}
\end{threeparttable}
\end{table}

% ══════════════════════════════════════════════════════════════════════════════
\section{Conclusion}
% ══════════════════════════════════════════════════════════════════════════════

This analysis examines the relationship between economic activity and CO2 emissions using cross-sectional OLS, first difference and fixed effects models on a panel of 192 countries between 1992 and 2024. The estimated elasticities declines substantially as the analysis move from cross-sectional OLS to estimators that absorb time-invariant confounders. This suggest that much of the raw GDP-CO2 correlation reflects between-country differences that remain persistent over time rather than a causal within-country effect. Adding urbanization rate as a confounder also modestly dampens the GDP coefficient, confirming a mild upward bias in the baseline estimates, though GDP remains a significant predictor of emissions across all specifications.

% ══════════════════════════════════════════════════════════════════════════════
\end{document}
% ══════════════════════════════════════════════════════════════════════════════